\dish{辣椒油}
\hiragana{làjiāoyóu}
\altdish{Sichuan chili oil}
\makes{250 mL}
\prep{30 minutes + 24--72 hours to mature}
\source{openai-chatgpt}
\label{laziyou}

\begin{ingredients}
  \ingr{250}{mL}{neutral oil (rapeseed or peanut)}
  \ingr{40--50}{g}{Sichuan chili flakes (ideally 70\% Erjingtiao / 30\% Chaotianjiao)}
  \ingr{1}{tsp}{sesame seeds}
  \ingrS{}{}{pinch salt}
  \ingr{1}{}{star anise}
  \ingr{1}{}{small cassia stick}
  \ingr{2}{}{bay leaves}
  \ingr{1}{tsp}{Sichuan pepper (optional)}
  \ingr{}{}{small slice ginger}
  \ingrS{2}{cloves}{garlic, lightly crushed}
\end{ingredients}


\begin{recipe}
  \begin{enumerate}

  \item Combine chili flakes, sesame seeds, and salt in a heatproof
    bowl.

  \item Heat the oil gently with the ginger, garlic, star anise,
    cassia, bay leaves, and Sichuan pepper. Hold at
    120--140\degreeC{}, 15--20~minutes; nothing should brown. Strain
    out the aromatics.

  \item Cool the oil to about 170\degreeC.

  \item Pour over the chili flakes in three additions. First third at
    170\degreeC{} (roasted aroma); wait 20--30~seconds. Second third
    at 140\degreeC{} (extracts color); wait 30~seconds. Final third at
    110\degreeC{} (preserves fresh chili fragrance). Stir.

  \item Cool completely before covering. Rest 24--72~hours before
    using.

  \end{enumerate}
\end{recipe}

\notes{} Keeps about a month at room temperature, several months
refrigerated. Use a clean spoon.

%\accord{}
